\section[Pair function and loss function at the edge]{The pair function and the loss function at the kinematic edge}\label{app:D.39}

This section collects the algebra behind \cref{sec:15.1,sec:15.2}: the closed form of the pair function $F_\mu$ and the bubbles of \cref{eq:15.1-bubbles}, the single quadrature formula for the conserving response \cref{eq:15.1-Rcons}, the loss-function ratio of \cref{eq:15.2-LBCS,eq:15.2-cons}, the evaluation of the Kramers--Kronig shift $\mathfrak{s}_{\mathrm{BCS}}$, and the sum rule that leaves no loss weight below the gap. The hypotheses are those of \cref{sec:15.1-closed}: $T = 0$, a constant density of states $\mathcal{N}_{\mathrm{F}}$ per spin orientation, a linearized band with $q \ll k_{\mathrm{F}}$, and no impurities. With the ``1'' of $\varepsilon$ dropped, $\varepsilon = -(k_{\mathrm{TF}}^2/q^2)R/2\mathcal{N}_{\mathrm{F}}$, and every ratio below is a function of $y = \omega/\Delta_{\mathrm{BCS}}$ and $qv_{\mathrm{F}}/\Delta_{\mathrm{BCS}}$.

\subsection*{The pair function}

The combination $(E + E')^2 - v_{\mathrm{F}}^2q^2\mu^2$ of \cref{eq:15.1-ident} is never smaller than $4\Delta_{\mathrm{BCS}}^2$. Write $\xi = \Delta_{\mathrm{BCS}}\sinh a$ and $\xi' = \Delta_{\mathrm{BCS}}\sinh b$, so that $E + E' = 2\Delta_{\mathrm{BCS}}\cosh A\cosh B$ and $\xi' - \xi = v_{\mathrm{F}}q\mu = 2\Delta_{\mathrm{BCS}}\cosh A\sinh B$ with $A = (a + b)/2$ and $B = (b - a)/2$. Then $(E + E')^2 - v_{\mathrm{F}}^2q^2\mu^2 = 4\Delta_{\mathrm{BCS}}^2\cosh^2A$. At fixed $\mu$, that is at fixed $\xi' - \xi$, the relation $\cosh b\,\dd b = \cosh a\,\dd a$ gives $\dd A = \dd a\,(\cosh a + \cosh b)/2\cosh b$; with $\dd\xi = \Delta_{\mathrm{BCS}}\cosh a\,\dd a$ the measure becomes $\dd\xi\,(E + E')/EE' = 2\,\dd A$. The same substitution gives $\int\dd\xi\,(E + E')/\{EE'[(E + E')^2 - v_{\mathrm{F}}^2q^2\mu^2]\} = \int\dd A/(2\Delta_{\mathrm{BCS}}^2\cosh^2A) = 1/\Delta_{\mathrm{BCS}}^2$ for every $\mu$. For the pair function of \cref{eq:15.1-F} it gives $F_\mu = (\mathcal{N}_{\mathrm{F}}/4\Delta_{\mathrm{BCS}}^2)\int\dd A/(\mathfrak{p}_\mu - \cosh^2A)$, and with $t = \tanh A$ the integral becomes $\int_{-1}^1\dd t/(\mathfrak{p}_\mu - 1 - \mathfrak{p}_\mu t^2)$, which is elementary. It gives \cref{eq:15.1-F} for $0 < \mathfrak{p}_\mu < 1$, and for the other ranges
\begin{equation}
  \frac{\arcsin\sqrt{\mathfrak{p}}}{\sqrt{\mathfrak{p}(1 - \mathfrak{p})}} \to
  \frac{\operatorname{arsinh}\sqrt{|\mathfrak{p}|}}{\sqrt{|\mathfrak{p}|(1 + |\mathfrak{p}|)}}\quad(\mathfrak{p} < 0), \qquad
  \to -\frac{\operatorname{arcosh}\sqrt{\mathfrak{p}} - \ii\pi/2}{\sqrt{\mathfrak{p}(\mathfrak{p} - 1)}}\quad(\mathfrak{p} > 1) ,
  \label{eq:D.39-Fcont}
\end{equation}
where the sign of the imaginary part is that of the retarded function, $\omega + \ii0^+$.

The density bubble has the integrand $|M^{\mathrm{N}}_{\boldsymbol{\eta}}|^2\cdot 2(E + E')/[\omega^2 - (E + E')^2]$. With $W = E + E'$ the partial fractions
\begin{equation*}
  \frac{W^2}{(W^2 - v_{\mathrm{F}}^2q^2\mu^2)(\omega^2 - W^2)} = \frac{1}{\omega_\mu^2}\Bigl[\frac{\omega^2}{\omega^2 - W^2} + \frac{v_{\mathrm{F}}^2q^2\mu^2}{W^2 - v_{\mathrm{F}}^2q^2\mu^2}\Bigr] ,
\end{equation*}
split it into an $F_\mu$ term and a term that the integral identity above evaluates; this is $R^0_{\boldsymbol{\eta}\boldsymbol{\eta}}$ of \cref{eq:15.1-bubbles}. The mixed and phase bubbles follow from the products of \cref{sec:15.1-closed} with $F_\mu$ alone: $R^0_{\mathsf{s}_2\boldsymbol{\eta}} = 2\ii\Delta_{\mathrm{BCS}}\omega\langle F_\mu\rangle_\mu$, and $R^0_{\mathsf{s}_2\mathsf{s}_2}(q,\omega) - R^0_{\mathsf{s}_2\mathsf{s}_2}(0,0) = \langle\omega_\mu^2F_\mu\rangle_\mu$. The logarithmic dependence of $R^0_{\mathsf{s}_2\mathsf{s}_2}$ on the cutoff $\omega_{\mathrm{D}}$ cancels in that difference: the cutoff-dependent part is $\mathcal{N}_{\mathrm{F}}\int\dd\xi/E$ over the shell $|\xi| < \omega_{\mathrm{D}}$, and the Doppler shift $\xi' - \xi = v_{\mathrm{F}}q\mu$ moves its limits by $v_{\mathrm{F}}q\mu$, a change of relative order $v_{\mathrm{F}}q/\omega_{\mathrm{D}}$. The gap equation fixes $R^0_{\mathsf{s}_2\mathsf{s}_2}(0,0) = -2/V_{\mathrm{BCS}}$, which gives the last of \cref{eq:15.1-bubbles}.

\subsection*{Quadrature of the conserving response}

The factor $F_\mu$ has an integrable inverse-square-root edge where $\omega_\mu = 2\Delta_{\mathrm{BCS}}$, at $v_{\mathrm{F}}q\mu = s_* = (\omega^2 - 4\Delta_{\mathrm{BCS}}^2)^{1/2}$. The substitution $v_{\mathrm{F}}q\mu = s_* \mp t^2$ on the two sides removes it, after which adaptive quadrature of the three averages $\langle F_\mu\rangle_\mu$, $\langle\omega_\mu^2F_\mu\rangle_\mu$ and the bubble of \cref{eq:15.1-Rcons} converges to $10^{-11}$. This one quadrature gives every value of $R^{\mathrm{cons}}$, $\mathfrak{v}$ and $\delta$ in \cref{sec:15.1,sec:15.2}. With $R^0$ the Lindhard bubble, $\mathfrak{g} = \Imag R^{\mathrm{cons}}/\Imag R^0 = \mathfrak{g}_{\mathrm{I}}(1 + \mathfrak{v})$, since $\Imag R^0_{\boldsymbol{\eta}\boldsymbol{\eta}} = \mathfrak{g}_{\mathrm{I}}\Imag R^0$, and $1 + \delta = \Real R^{\mathrm{cons}}/\Real R^0$, the quantities of \cref{eq:15.2-cons}. On the kinematic edge at $y = 3$ this gives $\mathfrak{v} = 0.214$ and $1 + \delta = 0.596$; the bare bubble computed in the same way has $1 + \delta = 1.329$, against $1 + \delta_{\mathfrak{s}} = 1.318$ from $\mathfrak{s}_{\mathrm{BCS}}$.

\subsection*{The loss-function ratio}

Let a response have $\Imag\varepsilon = \mathfrak{g}\,\Imag\varepsilon_{\mathrm{Lin}}$ and $\Real\varepsilon = (1 + \delta)\Real\varepsilon_{\mathrm{Lin}}$ at some $(q,\omega)$, with $\mathfrak{g}$ and $\delta$ real. Dividing $L = \Imag\varepsilon/[(\Real\varepsilon)^2 + (\Imag\varepsilon)^2]$ by the same expression for $\varepsilon_{\mathrm{Lin}}$, and then numerator and denominator by $(\Real\varepsilon_{\mathrm{Lin}})^2$,
\begin{equation}
  \frac{L}{L^{\mathrm{nor}}} = \frac{\mathfrak{g}\,\Imag\varepsilon_{\mathrm{Lin}}}{(1 + \delta)^2(\Real\varepsilon_{\mathrm{Lin}})^2 + \mathfrak{g}^2(\Imag\varepsilon_{\mathrm{Lin}})^2}\cdot\frac{(\Real\varepsilon_{\mathrm{Lin}})^2 + (\Imag\varepsilon_{\mathrm{Lin}})^2}{\Imag\varepsilon_{\mathrm{Lin}}}
  = \mathfrak{g}\,\frac{1 + \mathfrak{r}_\varepsilon^2}{(1 + \delta)^2 + \mathfrak{g}^2\mathfrak{r}_\varepsilon^2} ,
  \label{eq:D.39-resum}
\end{equation}
with $\mathfrak{r}_\varepsilon = \Imag\varepsilon_{\mathrm{Lin}}/\Real\varepsilon_{\mathrm{Lin}} = \pi\mathfrak{w}/2\Real\mathfrak{f}_{\mathrm{L}}(\mathfrak{w})$ from \cref{eq:15.2-Lin}. No approximation enters. The ratio is positive, and for $\delta \ge 0$ it is at most $\mathfrak{g}(1 + \mathfrak{r}_\varepsilon^2)/(1 + \mathfrak{g}^2\mathfrak{r}_\varepsilon^2)$.

The bare bubble has $\mathfrak{g} = \mathfrak{g}_{\mathrm{I}}(y)$ for $\omega < qv_{\mathrm{F}}$ by \cref{eq:15.1-gI}. Its shift follows from \cref{eq:15.1-sBCS}, $\Real\varepsilon^{\mathrm{bare}} - \Real\varepsilon_{\mathrm{Lin}} = (k_{\mathrm{TF}}^2/q^2)(\Delta_{\mathrm{BCS}}/qv_{\mathrm{F}})\mathfrak{s}_{\mathrm{BCS}}$; dividing by $\Real\varepsilon_{\mathrm{Lin}} = (k_{\mathrm{TF}}^2/q^2)\Real\mathfrak{f}_{\mathrm{L}}(\mathfrak{w})$ gives $\delta = \delta_{\mathfrak{s}}$ of \cref{eq:15.2-LBCS}.

\subsection*{The Kramers--Kronig shift}

The function $\mathfrak{s}_{\mathrm{BCS}}(y) = \operatorname{PV}\!\int_0^\infty\dd y'\,y'^2\delta\mathfrak{g}(y')/(y'^2 - y^2)$, with $\delta\mathfrak{g} = \mathfrak{g}_{\mathrm{I}}\Theta(y' - 2) - 1$, is evaluated as follows. With $\mathfrak{h}(y') = y'^2\delta\mathfrak{g}(y')/(y' + y)$ the integrand is $\mathfrak{h}(y')/(y' - y)$, and for an upper limit $Y > y$
\begin{equation*}
  \operatorname{PV}\!\int_0^Y\dd y'\,\frac{\mathfrak{h}(y')}{y' - y} = \int_0^Y\dd y'\,\frac{\mathfrak{h}(y') - \mathfrak{h}(y)}{y' - y} + \mathfrak{h}(y)\ln\frac{Y - y}{y} ,
\end{equation*}
whose first integrand is regular at $y' = y$. For $Y \to \infty$ the large-$y'$ form $\mathfrak{g}_{\mathrm{I}} - 1 = [2\ln(2y') - 1]/y'^2 + O(\ln y'/y'^4)$ makes the integrand of the original integral $[2\ln(2y') - 1]/y'^2$, whose integral from $Y_{\mathrm{tail}}$ to infinity is $[2\ln(2Y_{\mathrm{tail}}) + 1]/Y_{\mathrm{tail}}$; we integrate numerically to $Y_{\mathrm{tail}} = 10^4$ and add this tail. Near $y = 2$ the function $\mathfrak{h}$ steps by $(\pi/2)\cdot4/(2 + y) \to \pi/2$ at $y' = 2$, and a step $J$ in $\mathfrak{h}$ at $y' = 2$ adds $J\ln[1/|y - 2|]$ to the principal value; hence $\mathfrak{s}_{\mathrm{BCS}} \simeq (\pi/2)\ln[1/|y - 2|] + O(1)$. With $\delta_{\mathfrak{s}} \propto \ln[1/(y - 2)]$, \cref{eq:D.39-resum} falls as $(1 + \delta_{\mathfrak{s}})^{-2}$ toward the threshold once $\delta_{\mathfrak{s}}$ exceeds one.

\subsection*{No loss weight below the gap}

Let $\mathfrak{k}(\omega) = 2\mathcal{N}_{\mathrm{F}}/R(q,\omega)$, which is analytic in the upper half plane where $R$ has no zero, with $\Imag\mathfrak{k} \ge 0$ on the real axis because $\Imag R \le 0$. For the normal metal the moments \cref{eq:4.5-moments} at $q \ll k_{\mathrm{F}}$ give $R^0 = (\nd_eq^2/m_e\omega^2)[1 + \frac35(qv_{\mathrm{F}}/\omega)^2 + \dots]$, so $\mathfrak{k} = 3\omega^2/(qv_{\mathrm{F}})^2 - \frac95 + O(\omega^{-2})$, and $\mathfrak{k}(0) = -1$ by the static limit. The function $\tilde{\mathfrak{k}} = \mathfrak{k} - 3\omega^2/(qv_{\mathrm{F}})^2$ tends to a constant, and the half-range Kramers--Kronig relation \cref{eq:4.3-KKhalf} for $\tilde{\mathfrak{k}}$ at $\omega = 0$ reads
\begin{equation*}
  \frac{2}{\pi}\int_0^\infty\dd\omega\,\frac{\Imag\mathfrak{k}(\omega)}{\omega} = \tilde{\mathfrak{k}}(0) - \tilde{\mathfrak{k}}(\infty) = -1 + \frac95 = \frac45 .
\end{equation*}
A zero of $R$ on the real axis below $2\Delta_{\mathrm{BCS}}$, which is a zero of $\varepsilon$ and a $\delta$-function in $L$, would contribute to the left side. The conserving response has the same two limits: numerically $\tilde{\mathfrak{k}} = -1.8000$ at $\omega = 10^3\Delta_{\mathrm{BCS}}$ for $qv_{\mathrm{F}} = 5$ and $12\,\Delta_{\mathrm{BCS}}$, and $\mathfrak{k}(0) = -1$ by \cref{sec:15.1-closed}. Its integral over $\omega > 2\Delta_{\mathrm{BCS}}$ alone is 0.8000 at both momenta, as is that of the normal metal. The Anderson--Bogoliubov pole of $R^{\mathrm{cons}}$ is a zero of $\mathfrak{k}$ and does not enter.
